\documentclass[10pt]{article}
\usepackage[margin=20mm]{geometry}
\usepackage{amsmath,amssymb,mathtools,fontspec,unicode-math,xurl,hyperref}
\IfFontExistsTF{DejaVu Serif}{\setmainfont{DejaVu Serif}}{\setmainfont{Latin Modern Roman}}
\setmathfont{latinmodern-math.otf}
\setlength{\parindent}{0pt}\setlength{\parskip}{6pt}\sloppy
\title{Technical verification of
Geometric G-functions and Atypicality}
\author{AI-assisted technical review}
\date{}
\begin{document}
\maketitle
Abstract: This report examines the acquired third arXiv version, concentrating on the period construction, algebraic relations, height bounds and their stated dependencies. A scalar-matrix obstruction invalidates the homogeneous-tensor description of the period torsor used on a core proof path. A direct monodromy argument repairs the particular orbit-containment step needed for nontriviality of relations. Two further local defects concern a Jacobian-rank inference and the asserted simplicity of an orthogonal group. An additional quotient-incidence step remains unresolved. These conclusions concern specific statements and proofs; they do not establish that the main finiteness theorems are false. Scientific adjudication remains pending.\par
\section*{Identity, version and review scope}
The reviewed work is David Urbanik, Geometric G-functions and Atypicality, arXiv:2301.01857v3, submitted on 26 July 2025, with printed title-page date 29 July 2025 [U]. Its frozen journal identity is Duke Mathematical Journal 174 (2025), 2425–2512, DOI https://doi.org/10.1215/00127094-2024-0079. The acquired author PDF has 72 pages; the journal pagination has 88 pages. The publisher’s full text was not acquired or certified equivalent. All page references below mean one-based pages of the acquired PDF; its printed Arabic page numbers agree.\par
The review followed the local log de Rham construction and realization in Sections 2–4; the algebra of relations and nontriviality argument in Sections 5–6; the point-counting reduction and idempotent bounds in Section 7; the dependence of the applications in Section 8; and the alternative comparison argument in Appendix B. The coverage ledger distinguishes direct inspection, bounded native checks and unverified external inputs. It does not assert a formal verification of every cited theorem. The released LaTeX was inspected as source text. No author program or author LaTeX was executed. A separate supplement or computational repository was not identified in the acquired materials; this absence is not a mathematical defect.\par
\section*{Assumptions relevant to the conclusions}
The period construction uses a projective family over a smooth curve, a simple normal-crossing degeneration, a local equation \(s=z_1\cdots z_\mu\), and cohomological degree \(w=\mu-1\). The height arguments further require the specified simplicity and CM conditions, two independent functionals where stated, relations at all relevant places, and functional nontriviality in the Hasse principle. CM is read with the paper’s algebraic Definition 1.4.\par
In Section 6 the connected algebraic monodromy group \(H\) is obtained as the identity component of the Zariski closure of monodromy, after the stated covering adjustment. The paper explicitly sets \[L^{a,b}=L^{\otimes a}\otimes(L^*)^{\otimes b}.\] Preserving a subspace setwise is distinct from fixing its vectors pointwise. This distinction is decisive for Problem 1.\par
\section*{Established substantive findings}
Problem 1 (A homogeneous-tensor stabilizer cannot give the claimed torsor).  The universal stabilizer assertion in Section 6.1, p. 38, is false under the displayed definition of \(L^{a,b}\). Consequently the stated justification of the \(H\)-torsor in equation (26), pp. 38–39, does not establish the construction used in Proposition 6.4.\par
\subsection*{Claim, parameters and complete check}
The source asserts that for every algebraic subgroup \(H\subset\operatorname{GL}(L)\) there are \(a,b\geq0\) and a line \(M\subset L^{a,b}\) whose setwise stabilizer is \(H\). Take \(L=\mathbb C^2\) and \(H=\operatorname{SL}_2\), which meet these hypotheses. For every \(a,b\) and every subspace \(M\), \[(\lambda I)|_{L^{a,b}}=\lambda^{a-b}I,\qquad
 (\lambda I)M=M\quad(\lambda\in\mathbb C^\times).\] Thus its stabilizer contains all scalar matrices. But \(\det(2I_2)=4\), so \(2I_2\notin\operatorname{SL}_2\). No choice permitted by the assertion can produce \(\operatorname{SL}_2\). Allowing higher-dimensional \(M\), as the following source paragraph does, cannot remove this obstruction.\par
The obstruction occurs in the actual local-system setting. On \(T=\mathbb P^1_\mathbb C\setminus\{0,1,\infty\}\) assign its two free fundamental-group generators to \[U=\begin{pmatrix}1&1\\0&1\end{pmatrix},\qquad
 V=\begin{pmatrix}1&0\\1&1\end{pmatrix}.\] The Zariski closures of their integral powers contain the upper and lower one-parameter unipotent groups: a polynomial vanishing on all integers vanishes identically. These groups generate \(\operatorname{SL}_2\). The resulting rank-two local system has connected algebraic monodromy \(\operatorname{SL}_2\) and fits the regular-singular Riemann–Hilbert setup invoked in Section 6.1.\par
For illustration, the natural determinant line \(M=\bigwedge^2L\subset L^{\otimes2}\) has setwise stabilizer \(\operatorname{GL}_2\). Putting that line into equation (26) gives fibre \(\operatorname{GL}_2\) at the basepoint, whereas parallel transport fixes a determinant tensor and lies in the determinant-one subvariety. The point \(2I_2\) is in the former fibre and outside the closure of the latter graph. This illustration is not presented as a line satisfying the impossible original assertion; it shows the consequence of substituting the natural invariant line into the displayed formula.\par
\subsection*{Counterevidence, dependence and impact}
The PDF confirms the single tensor bidegree and the setwise equation; this is not an extraction error. Material counterevidence appears in Section 7.4, p. 58: that paragraph allows a direct sum over tensor bidegrees and invokes pointwise-fixed tensors. Such a representation can distinguish scalars. For example, a line generated by \(1+t_1+\cdots+t_r\), with the tensors placed in separate summands, has setwise preservation forcing the coefficient of the trivial summand, and hence the individual tensors, to be fixed.\par
This later adjustment supplies a credible repair route, but does not make the explicit assertion on p. 38 true. The displayed dependence is Section 6.1 \(\longrightarrow\) Proposition 6.4 \(\longrightarrow\) Lemma 6.7 \(\longrightarrow\) Propositions 6.8–6.9 \(\longrightarrow\) Sections 6.5.2 and 6.6.2 \(\longrightarrow\) nontriviality in the Hasse-principle application. This is a demonstrated invalid derivation on a core proof path and supports Major under the supplied calibration. It is not a counterexample to Theorems 1.5 or 1.15. The separately adjusted Section 7.4 torsor has not been refuted.\par
\subsection*{A verified repair of the needed orbit-containment step}
The particular conclusion of Lemma 6.7 can be proved without that torsor. Let \(\widetilde R\) be the closed algebraic subset in Definition 6.6 containing the local flat graph \(\Gamma_1\). Continue the defining identities along each loop based at \(c\). They still hold, so \(\widetilde R_c\) contains the monodromy orbit of \(\alpha_1^*\). As \(\widetilde R_c\) is Zariski closed, it contains \(H(\mathbb C)\alpha_1^*\). With \(\widetilde R_c=V(Q)\) this is the required conclusion.\par
For a nonzero product of linear forms, irreducibility of a connected group orbit implies that one factor vanishes on the whole orbit. Its linear span is then a proper invariant subspace. When \(\alpha_1^*\) is \(E\)-defined, that span descends from the span of its \(E\)-monodromy orbit, contradicting the stipulated \(E\)-simplicity. This checks the specific nontriviality link, not every other input of the full height theorem.\par
Problem 2 (Uniqueness of a formal solution does not force the stated Jacobian rank).  In the proof of Lemma 4.3, the first paragraph of p. 32 infers full column rank of the matrix \(\boldsymbol M\) in equations (16)–(18) from uniqueness of the formal power-series solution. That inference is invalid.\par
\subsection*{Complete counterexample}
Set \(K=\mathbb Q\), \(\nu=\sigma=c=1\), \(q=0\), \(B_1(y)=y^3\) and \(C_1(x)=x^3\). These are the polynomials and initial condition permitted by Lemma 4.3 (p. 30). There is exactly one \(A\in\mathbb Q[[x]]\) with \(A(0)=0\) and \(A^3=x^3\). Indeed \(\operatorname{ord}_x(A)=1\), so \(A=xu\) and \(u^3=1\). Its rational constant term is \(1\), while \(u^2+u+1\) is a unit with constant term \(3\). Factoring \(u^3-1\) gives \(u=1\) and \(A=x\). Nevertheless, \[\boldsymbol M=(B_1'(0))=(0),\qquad
 \operatorname{rank}\boldsymbol M=0<1=\sigma .\] Later coefficient equations can constrain earlier apparent choices; a deficient matrix at each immediate recursion does not create infinitely many compatible full solutions.\par
\subsection*{Boundary and restricted repair}
The example is polynomial and does not refute the lemma’s convergence conclusion. The actual application in the preceding Lemma 4.2 has the stronger hypothesis that \(g:U\to V\) is étale. If a tangent vector annihilates the gradients of the equations defining \(U\), it lies in \(T_qU\); if it also annihilates \(dg\), étaleness makes it zero. Hence the combined Jacobian of the defining equations and \(g\) has full column rank independently of uniqueness. Using this rank fact repairs the recursion in that application. It does not prove the broader uniqueness-only lemma. The supported assessment is Minor, because the inspected core application has this independent justification and no material failure of the main realization theorem was demonstrated.\par
Problem 3 (The orthogonal rank-four case contradicts Lemma 7.21).  Lemma 7.21, p. 56, asserts that for a non-degenerate polarized Hodge structure the special isometry group is the \(\mathbb Q\)-simple special Mumford–Tate group of a point in the orbit. Its proof excludes only \(\operatorname{SO}_2\) as a possible non-simple case. A split \(\operatorname{SO}_4\) example meets the stated non-degeneracy hypothesis.\par
\subsection*{Construction and verification}
Let \(U_\mathbb Z=\mathbb Z^2\) with \(J=\left(\begin{smallmatrix}0&-1\\1&0\end{smallmatrix}\right)\), the weight-one Hodge structure whose \((1,0)\) part is the \(+i\) eigenspace of \(J\), and alternating polarization \(J\). Put \(W_\mathbb Z=U_\mathbb Z\otimes U_\mathbb Z\) with its tensor Hodge structure. Its Weil operator and symmetric polarization can be written \[C=J\otimes J,\qquad Q=J\otimes J,\qquad C^tQ=I_4.\] Thus \(Q(C-, -)\) is positive definite, while \(Q\) has signature \((2,2)\). The Hodge numbers are \((1,2,1)\), so \(\dim_\mathbb CF^1=3<4\); Definition 7.20 on p. 55 calls this non-degenerate. It is also realized by the Künneth summand \(H^1(E)\otimes H^1(E)\) of a product of elliptic curves.\par
The map \[\operatorname{SL}_2\times\operatorname{SL}_2\longrightarrow\operatorname{SO}(W,Q),\qquad
 (A,B)\longmapsto A\otimes B\] preserves \(Q\), has kernel \(\{(I,I),(-I,-I)\}\), and has six-dimensional image, the connected split special orthogonal group. The image of \(\operatorname{SL}_2\times\{I\}\) is a three-dimensional connected \(\mathbb Q\)-normal subgroup of this six-dimensional group. Consequently \(\operatorname{SO}(W,Q)\) is not \(\mathbb Q\)-simple. Changing the Hodge point cannot change this property of the fixed group.\par
\subsection*{Impact and repair boundary}
This refutes the simplicity clause and the proof’s classification. It does not by itself refute existence of generic full \(\operatorname{SO}(2,2)\) special Mumford–Tate group, nor Lemma 7.22(iii). The downstream abelian application in Proposition 7.27 expressly uses weight one, with symplectic factors. No core dependence on the false orthogonal simplicity assertion was demonstrated; the assessment is Minor. Restricting the invoked lemma to the weight-one setting avoids this example. Merely deleting an adjective has not been verified as a complete all-weight replacement.\par
\section*{Unresolved point and checks that did not establish defects}
In the proof of Corollary 7.17, p. 54, let \(A_j=\{g:gP_j\in I\}\). The source replaces the condition that \(gM(\mathbb C)P_j\) meet \(I\) by the condition \(gP_j\in I\). The former is invariant under right multiplication by \(M\); the latter need not be. The set-theoretically correct expression is \[q^{-1}(G_j)=A_jM(\mathbb C).\] Whenever a positive-dimensional orbit meets \(I\) in an isolated point, moving its representative slightly inside the orbit exhibits the distinction. Zero-dimensional orbits are counterevidence to a blanket objection. This review has not produced a complete geometric counterexample satisfying all global infinitude, height and Galois hypotheses of Corollary 7.17, nor verified a full replacement for the complex-analytic incidence argument preceding Lemma 7.18. This item is Needs Evidence; it is not used to raise the supported severity or to declare the corollary false.\par
Several initially suspicious steps have positive explanations. In Proposition 5.3 a largest cyclic block supplies a nonzero Krylov minor. For Proposition 5.4, if the characteristic and minimal polynomials of \(\tau\) agree, the nonzero polynomial \[D_\tau(x)=\det[x,\tau^*x,\ldots,(\tau^*)^{m-1}x]\] has degree \(m\), is defined over the base field, vanishes on noncyclic vectors, and factors into linear forms over a splitting field. It supplies a uniform relation for the relevant places. The two-functional argument in Section 6.6 uses the valid inequality \(\dim(E^2\cap W)\geq1\) for a codimension-one subspace \(W\).\par
For Lemma 7.25, the adjugate identity \[(\det\alpha)e=\operatorname{adj}(\alpha)e'\alpha\] shows that every coordinate denominator divides the isogeny degree when \(e'\) and \(\alpha\) are integral. Appendix A’s triangular support and symmetry argument yields \(\|s\eta s^{-1}\|_\infty\leq C D\), with \(D\) the largest coordinate denominator. Under the usual affine multiplicative height convention, one can use the common denominator to obtain \(H(s\eta s^{-1})\leq C'D^{m^2+1}\). This preserves the polynomial conclusion. The different possible matrix-height conventions are treated as a clarification, not as a claim that Proposition A.1 is false.\par
Finally, Appendix B, pp. 64–66, supplies an open-annulus decomposition and kernel description addressing the apparent splitting freedom in Section 3. The external comparison results remain dependencies rather than theorems independently re-proved in this review.\par
\section*{Overall dependence and state}
Problem 1 affects the displayed proof of a necessary nontriviality step in the height method. The direct monodromy repair above supports that step, while the mixed-tensor paragraph in Section 7.4 is essential counterevidence. Problems 2–3 concern stated auxiliary arguments beyond what their checked core uses require. No counterexample to the main height or finiteness conclusions has been established.\par
The arXiv version page and the author’s publication page checked for this review still identified v3; the bounded check did not identify a specific erratum resolving these findings. This is not an assertion of first discovery or absence of all later corrections. A source-by-source comparison with the publisher’s final text remains outside this report.\par
Item | Assessment | Repair verification | Adoption\par
Homogeneous torsor | Major | Direct orbit argument checked; full theorem not certified | Unknown\par
Jacobian inference | Minor | Étale application repaired; general lemma open here | Unknown\par
Orthogonal simplicity | Minor | Weight-one restriction supported | Unknown\par
Quotient incidence | Needs Evidence | Full replacement not verified | Unknown\par
Scientific decision is pending. The human-review record is empty. This initial review is not independent acceptance.\par
\section*{Supporting checks}
Analytic proofs above are distinguished from finite matrix checks and static source inspection. Imported inputs include the Hasse principle, external field-degree and isogeny bounds, functional transcendence and the comparison theorems. These were not all independently re-proved.\par
The primary executor integrated three native in-conversation bounded checks. They are AI checks, not human review or independent acceptance. The external controller’s visible model observation and the executor/helpers’ runtime self-description are separately recorded in provenance; backend equivalence is not asserted. No author outreach or adoption occurred.\par
9 D. Urbanik, Geometric G-functions and Atypicality, arXiv:2301.01857v3, 26 July 2025, acquired author PDF and released LaTeX source, https://arxiv.org/abs/2301.01857v3. D. Urbanik, publication list, https://www.math.ias.edu/\textasciitilde{}dburbanik/, consulted 26 October 2026 for version and correction context only.\par
\end{document}
